\setlength{\tabcolsep}{0pt}
\begin{longtable}{|C{12.5mm}|C{12.5mm}|C{12.5mm}|C{12.5mm}|C{12.5mm}|C{12.5mm}|C{12.5mm}|C{12.5mm}|C{12.5mm}|C{12.5mm}|}
\hline
{\bf{HEX}} & {\bf{DEC}} & {\bf{DB7}} & {\bf{DB6}} & {\bf{DB5}} & {\bf{DB4}} & {\bf{DB3}} & {\bf{DB2}} & {\bf{DB1}} & {\bf{DB0}} \\
\hline
\endfirsthead
\multicolumn{3}{l@{}}{\ldots vervolgd}\\
\hline
{\bf{HEX}} & {\bf{DEC}} & {\bf{DB7}} & {\bf{DB6}} & {\bf{DB5}} & {\bf{DB4}} & {\bf{DB3}} & {\bf{DB2}} & {\bf{DB1}} & {\bf{DB0}} \\
\endhead
\multicolumn{3}{l@{}}{vervolgd \ldots}\\
\endfoot
\hline
\endlastfoot
\footnotesize  D609\index{Registers!\$D609} \index{Registers!54793} & \footnotesize 54793 & \multicolumn{7}{c|}{\footnotesize --}& \footnotesize UFAST \\
\hline
\footnotesize  D60A\index{Registers!\$D60A} \index{Registers!54794} & \footnotesize 54794 & \footnotesize KEYQUEUE & \footnotesize MODKEYCAPS & \footnotesize MODKEYSCRL & \footnotesize MODKEYALT & \footnotesize MODKEYMEGA & \footnotesize MODKEYCTRL & \footnotesize MODKEYRSHFT & \footnotesize MODKEYLSHFT \\
\hline
\footnotesize  D60B\index{Registers!\$D60B} \index{Registers!54795} & \footnotesize 54795 & \footnotesize OSKZEN & \footnotesize OSKZON & \multicolumn{6}{c|}{\footnotesize PORTF\index{Registers!PORTF}}\\
\hline
\footnotesize  D60C\index{Registers!\$D60C} \index{Registers!54796} & \footnotesize 54796 & \multicolumn{2}{c|}{\footnotesize PORTFDDR\index{Registers!PORTFDDR}}& \multicolumn{6}{c|}{\footnotesize PORTFDDR\index{Registers!PORTFDDR}}\\
\hline
\footnotesize  D60D\index{Registers!\$D60D} \index{Registers!54797} & \footnotesize 54797 & \footnotesize HDSCL & \footnotesize HDSDA & \footnotesize SDBSH & \footnotesize SDCS & \footnotesize SDCLK & \footnotesize SDDATA & \footnotesize RST41 & \footnotesize CONN41 \\
\hline
\footnotesize  D60E\index{Registers!\$D60E} \index{Registers!54798} & \footnotesize 54798 & \multicolumn{8}{c|}{\footnotesize BASHDDR\index{Registers!BASHDDR}}\\
\hline
\footnotesize  D60F\index{Registers!\$D60F} \index{Registers!54799} & \footnotesize 54799 & \footnotesize ACCESSKEY & \footnotesize OSKDIM & \footnotesize REALHW & \multicolumn{3}{c|}{\footnotesize --}& \footnotesize KEYUP & \footnotesize KEYLEFT \\
\hline
\footnotesize  D610\index{Registers!\$D610} \index{Registers!54800} & \footnotesize 54800 & \multicolumn{8}{c|}{\footnotesize ASCIIKEY\index{Registers!ASCIIKEY}}\\
\hline
\footnotesize  D611\index{Registers!\$D611} \index{Registers!54801} & \footnotesize 54801 & \footnotesize MDISABLE & \footnotesize MCAPS & \footnotesize MSCRL & \footnotesize MALT & \footnotesize MMEGA & \footnotesize MCTRL & \footnotesize MRSHFT & \footnotesize MLSHFT \\
\hline
\footnotesize  D612\index{Registers!\$D612} \index{Registers!54802} & \footnotesize 54802 & \footnotesize LJOYB & \footnotesize LJOYA & \footnotesize JOYSWAP & \footnotesize OSKDEBUG & \multicolumn{4}{c|}{\footnotesize --}\\
\hline
\footnotesize  D613\index{Registers!\$D613} \index{Registers!54803} & \footnotesize 54803 & \multicolumn{8}{c|}{\footnotesize KEYMATRIXPEEK\index{Registers!KEYMATRIXPEEK}}\\
\hline
\footnotesize  D614\index{Registers!\$D614} \index{Registers!54804} & \footnotesize 54804 & \multicolumn{8}{c|}{\footnotesize KEYMATRIXSEL\index{Registers!KEYMATRIXSEL}}\\
\hline
\footnotesize  D615\index{Registers!\$D615} \index{Registers!54805} & \footnotesize 54805 & \footnotesize OSKEN & \multicolumn{7}{c|}{\footnotesize VIRTKEY1\index{Registers!VIRTKEY1}}\\
\hline
\footnotesize  D616\index{Registers!\$D616} \index{Registers!54806} & \footnotesize 54806 & \footnotesize OSKALT & \multicolumn{7}{c|}{\footnotesize VIRTKEY2\index{Registers!VIRTKEY2}}\\
\hline
\footnotesize  D617\index{Registers!\$D617} \index{Registers!54807} & \footnotesize 54807 & \footnotesize OSKTOP & \multicolumn{7}{c|}{\footnotesize VIRTKEY3\index{Registers!VIRTKEY3}}\\
\hline
\footnotesize  D618\index{Registers!\$D618} \index{Registers!54808} & \footnotesize 54808 & \multicolumn{8}{c|}{\footnotesize KSCNRATE\index{Registers!KSCNRATE}}\\
\hline
\footnotesize  D619\index{Registers!\$D619} \index{Registers!54809} & \footnotesize 54809 & \multicolumn{8}{c|}{\footnotesize PETSCIIKEY\index{Registers!PETSCIIKEY}}\\
\hline
\footnotesize  D61A\index{Registers!\$D61A} \index{Registers!54810} & \footnotesize 54810 & \multicolumn{8}{c|}{\footnotesize SYSCTL\index{Registers!SYSCTL}}\\
\hline
\footnotesize  D61D\index{Registers!\$D61D} \index{Registers!54813} & \footnotesize 54813 & \footnotesize KEYLED\-ENA & \multicolumn{7}{c|}{\footnotesize KEYLED\-REG\index{Registers!KEYLEDREG}}\\
\hline
\footnotesize  D61E\index{Registers!\$D61E} \index{Registers!54814} & \footnotesize 54814 & \multicolumn{8}{c|}{\footnotesize KEYLED\-VAL\index{Registers!KEYLEDVAL}}\\
\hline
\footnotesize  D620\index{Registers!\$D620} \index{Registers!54816} & \footnotesize 54816 & \multicolumn{8}{c|}{\footnotesize POTAX\index{Registers!POTAX}}\\
\hline
\footnotesize  D621\index{Registers!\$D621} \index{Registers!54817} & \footnotesize 54817 & \multicolumn{8}{c|}{\footnotesize POTAY\index{Registers!POTAY}}\\
\hline
\footnotesize  D622\index{Registers!\$D622} \index{Registers!54818} & \footnotesize 54818 & \multicolumn{8}{c|}{\footnotesize POTBX\index{Registers!POTBX}}\\
\hline
\footnotesize  D623\index{Registers!\$D623} \index{Registers!54819} & \footnotesize 54819 & \multicolumn{8}{c|}{\footnotesize POTBY\index{Registers!POTBY}}\\
\hline
\footnotesize  D625\index{Registers!\$D625} \index{Registers!54821} & \footnotesize 54821 & \multicolumn{8}{c|}{\footnotesize J21L\index{Registers!J21L}}\\
\hline
\footnotesize  D626\index{Registers!\$D626} \index{Registers!54822} & \footnotesize 54822 & \multicolumn{8}{c|}{\footnotesize J21H\index{Registers!J21H}}\\
\hline
\footnotesize  D627\index{Registers!\$D627} \index{Registers!54823} & \footnotesize 54823 & \multicolumn{8}{c|}{\footnotesize J21LDDR\index{Registers!J21LDDR}}\\
\hline
\footnotesize  D628\index{Registers!\$D628} \index{Registers!54824} & \footnotesize 54824 & \multicolumn{4}{c|}{\footnotesize BOARDMINOR\index{Registers!BOARDMINOR}}& \multicolumn{4}{c|}{\footnotesize J21HDDR\index{Registers!J21HDDR}}\\
\hline
\footnotesize  D629\index{Registers!\$D629} \index{Registers!54825} & \footnotesize 54825 & \multicolumn{8}{c|}{\footnotesize M65MODEL\index{Registers!M65MODEL}}\\
\hline
\end{longtable}
\begin{itemize}
\item{\bf{ACCESSKEY}\index{Registers!ACCESSKEY}} Toegankelijke toetsenbordinvoer inschakelen via de vuurknop van joystickpoort 2
\item{\bf{ASCIIKEY}\index{Registers!ASCIIKEY}} Bovenkant van de typgebeurteniswachtrij als ASCII. Schrijven om de gebeurtenis te wissen en de volgende gereed te maken.
\item{\bf{BASHDDR}\index{Registers!BASHDDR}} Data Direction Register (DDR) voor de \$D60D-bitbashingpoort.
\item{\bf{BOARDMINOR}\index{Registers!BOARDMINOR}} PCB-minorrevisie lezen (enkel R5+, anders worden nullen gelezen)
\item{\bf{CONN41}\index{Registers!CONN41}} Interne 1541-stationskoppeling (1=het interne 1541-station op de IEC-bus aansluiten)
\item{\bf{HDSCL}\index{Registers!HDSCL}} HDMI I2C-besturingsinterface SCL-klok
\item{\bf{HDSDA}\index{Registers!HDSDA}} HDMI I2C-besturingsinterface SDA-datalijn
\item{\bf{J21H}\index{Registers!J21H}} J21-pinnen 11 -- 14 invoer/uitvoerwaarden
\item{\bf{J21HDDR}\index{Registers!J21HDDR}} J21-pinnen 11 -- 14 data direction register
\item{\bf{J21L}\index{Registers!J21L}} J21-pinnen 1 -- 6, 9 -- 10 invoer/uitvoerwaarden
\item{\bf{J21LDDR}\index{Registers!J21LDDR}} J21-pinnen 1 -- 6, 9 -- 10 data direction register
\item{\bf{JOYSWAP}\index{Registers!JOYSWAP}} Joystickpoorten 1 \& 2 verwisselen
\item{\bf{KEYLEDENA}\index{Registers!KEYLEDENA}} Besturing van de toetsenbord-LED's inschakelen
\item{\bf{KEYLEDREG}\index{Registers!KEYLEDREG}} Registerkeuze voor de toetsenbord-LED's (R,G,B-kanalen x 4 = 0 tot 11)
\item{\bf{KEYLEDVAL}\index{Registers!KEYLEDVAL}} Waarde van het toetsenbord-LED-register (enkel schrijven)
\item{\bf{KEYLEFT}\index{Registers!KEYLEFT}} De C65 Cursor links-toets rechtstreeks lezen
\item{\bf{KEYMATRIXPEEK}\index{Registers!KEYMATRIXPEEK}} 8-bits segment van de gecombineerde toetsenbordmatrix (LEZEN)
\item{\bf{KEYMATRIXSEL}\index{Registers!KEYMATRIXSEL}} Bepaalt welk 8-bits segment van de gecombineerde toetsenbordmatrix gelezen wordt.
\item{\bf{KEYQUEUE}\index{Registers!KEYQUEUE}} 1 = De typgebeurteniswachtrij is niet leeg. Schrijf 0 naar deze bit om de wachtrij te spoelen.
\item{\bf{KEYUP}\index{Registers!KEYUP}} De C65 Cursor omhoog-toets rechtstreeks lezen
\item{\bf{KSCNRATE}\index{Registers!KSCNRATE}} Fysieke scanfrequentie van het toetsenbord (\$00=50MHz, \$FF=$\sim$200KHz)
\item{\bf{LJOYA}\index{Registers!LJOYA}} Roteert de invoer van joystick A met 180 graden (voor linkshandig gebruik)
\item{\bf{LJOYB}\index{Registers!LJOYB}} Roteert de invoer van joystick B met 180 graden (voor linkshandig gebruik)
\item{\bf{M65MODEL}\index{Registers!M65MODEL}} MEGA65-model-ID. Kan worden gebruikt om het model van de MEGA65 te bepalen waarop een programma wordt uitgevoerd, bijv. om aanraakbediening op de MEGAphone in te schakelen.
\item{\bf{MALT}\index{Registers!MALT}} Toestand van de ALT-toets (direct; alleen lezen).
\item{\bf{MCAPS}\index{Registers!MCAPS}} Toestand van de CAPS LOCK-toets (direct; alleen lezen).
\item{\bf{MCTRL}\index{Registers!MCTRL}} Toestand van de CTRL-toets (direct; alleen lezen).
\item{\bf{MDISABLE}\index{Registers!MDISABLE}} Schakelt de modificatietoetsen uit.
\item{\bf{MLSHFT}\index{Registers!MLSHFT}} Toestand van de linker SHIFT-toets (direct; alleen lezen).
\item{\bf{MMEGA}\index{Registers!MMEGA}} Toestand van de MEGA/C=-toets (direct; alleen lezen).
\item{\bf{MODKEYALT}\index{Registers!MODKEYALT}} Toestand van de ALT-toets bovenaan de wachtrij van typegebeurtenissen. 1 = ingedrukt tijdens de gebeurtenis.
\item{\bf{MODKEYCAPS}\index{Registers!MODKEYCAPS}} Toestand van de CAPS LOCK-toets bovenaan de wachtrij van typegebeurtenissen. 1 = ingedrukt tijdens de gebeurtenis.
\item{\bf{MODKEYCTRL}\index{Registers!MODKEYCTRL}} Toestand van de CTRL-toets bovenaan de wachtrij van typegebeurtenissen. 1 = ingedrukt tijdens de gebeurtenis.
\item{\bf{MODKEYLSHFT}\index{Registers!MODKEYLSHFT}} Toestand van de linker SHIFT-toets bovenaan de wachtrij van typegebeurtenissen. 1 = ingedrukt tijdens de gebeurtenis.
\item{\bf{MODKEYMEGA}\index{Registers!MODKEYMEGA}} Toestand van de MEGA/C=-toets bovenaan de wachtrij van typegebeurtenissen. 1 = ingedrukt tijdens de gebeurtenis.
\item{\bf{MODKEYRSHFT}\index{Registers!MODKEYRSHFT}} Toestand van de rechter SHIFT-toets bovenaan de wachtrij van typegebeurtenissen. 1 = ingedrukt tijdens de gebeurtenis.
\item{\bf{MODKEYSCRL}\index{Registers!MODKEYSCRL}} Toestand van de NOSCRL-toets bovenaan de wachtrij van typegebeurtenissen. 1 = ingedrukt tijdens de gebeurtenis.
\item{\bf{MRSHFT}\index{Registers!MRSHFT}} Toestand van de rechter SHIFT-toets (direct; alleen lezen).
\item{\bf{MSCRL}\index{Registers!MSCRL}} Toestand van de NOSCRL-toets (direct; alleen lezen).
\item{\bf{OSKALT}\index{Registers!OSKALT}} Toont de alternatieve indeling van het schermtoetsenbord (gewoonlijk het nummerpad van de MEGA65-telefoon)
\item{\bf{OSKDEBUG}\index{Registers!OSKDEBUG}} Debugt de OSK-overlay (alleen schrijven)
\item{\bf{OSKDIM}\index{Registers!OSKDIM}} Dimt het achtergrondmateriaal achter het schermtoetsenbord licht of sterk
\item{\bf{OSKEN}\index{Registers!OSKEN}} Schakelt de weergave van de samengestelde overlay van het schermtoetsenbord in
\item{\bf{OSKTOP}\index{Registers!OSKTOP}} 1=toont het schermtoetsenbord bovenaan, 0=toont het schermtoetsenbord onderaan het scherm.
\item{\bf{OSKZEN}\index{Registers!OSKZEN}} Toont de hardwarezoom van het gebied onder het eerste aanraakpunt voor het schermtoetsenbord
\item{\bf{OSKZON}\index{Registers!OSKZON}} Toont altijd de hardwarezoom van het gebied onder het eerste aanraakpunt
\item{\bf{PETSCIIKEY}\index{Registers!PETSCIIKEY}} Bovenaan de wachtrij van typegebeurtenissen als PETSCII. Schrijf erin om de gebeurtenis te wissen en gereed te maken voor de volgende.
\item{\bf{PORTF}\index{Registers!PORTF}} PMOD-poort A op het FPGA-bord (gegevens) (alleen Nexys4-borden)
\item{\bf{PORTFDDR}\index{Registers!PORTFDDR}} PMOD-poort A op het FPGA-bord (DDR)
\item{\bf{POTAX}\index{Registers!POTAX}} Leest de paddle X van poort A, zonder dat u met de SID/CIA-instellingen hoeft te rommelen.
\item{\bf{POTAY}\index{Registers!POTAY}} Leest de paddle Y van poort A, zonder dat u met de SID/CIA-instellingen hoeft te rommelen.
\item{\bf{POTBX}\index{Registers!POTBX}} Leest de paddle X van poort B, zonder dat u met de SID/CIA-instellingen hoeft te rommelen.
\item{\bf{POTBY}\index{Registers!POTBY}} Leest de paddle Y van poort B, zonder dat u met de SID/CIA-instellingen hoeft te rommelen.
\item{\bf{REALHW}\index{Registers!REALHW}} Zet op 1 indien de MEGA65 op echte hardware wordt uitgevoerd, zet op 0 indien geëmuleerd (Xemu) of gesimuleerd (ghdl).
\item{\bf{RST41}\index{Registers!RST41}} Interne reset van het 1541-schijfstation (1=reset, 0=operationeel)
\item{\bf{SDBSH}\index{Registers!SDBSH}} Bitbash-modus van de SD-kaart inschakelen
\item{\bf{SDCLK}\index{Registers!SDCLK}} SD-kaart SCLK
\item{\bf{SDCS}\index{Registers!SDCS}} SD-kaart CS\_BO
\item{\bf{SDDATA}\index{Registers!SDDATA}} SD-kaart MOSI/MISO
\item{\bf{SYSCTL}\index{Registers!SYSCTL}} Systeembeheervlagen (doelapparaatspecifiek)
\item{\bf{UFAST}\index{Registers!UFAST}} UART-baudklokbron van de C65: 1 = 7,09375 MHz, 0 = 80 MHz (VIC-IV-pixelpitch)
\item{\bf{VIRTKEY1}\index{Registers!VIRTKEY1}} Op \$7F zetten voor geen toetsaanslag; geef anders het virtuele toetsaanslag aan
\item{\bf{VIRTKEY2}\index{Registers!VIRTKEY2}} Op \$7F zetten voor geen toetsaanslag; geef anders de 2e virtuele toetsaanslag aan
\item{\bf{VIRTKEY3}\index{Registers!VIRTKEY3}} Op \$7F zetten voor geen toetsaanslag; geef anders de 3e virtuele toetsaanslag aan
\end{itemize}
